\section{Benchmarking and competitive evaluation}
\label{sec:benchmarking}
\label{app:benchmarking}

\paragraph{Primary architecture and historical evidence.}
Primary Prometheus is the latest fully fused architecture actually validated
under TODO~3, with TODO~2's independent benefit also validated before the
Primary freeze. Its core mechanism, code version, and evaluation configuration
are frozen before actual retraining across all five benchmark layers. The
earlier native/coordinate-renderer streams are historical results and internal
ablations only; their outcomes, parameters, and costs are not pooled with or
substituted for new Primary measurements. The limited D3 hybrid V2 diagnostic
is development evidence, not completed TODO~3 validation or a Primary freeze.

\paragraph{Baseline validity is independent of accuracy.}
A configuration is valid evidence when its implementation is checked, any
claimed reproduction is faithful, information is fair, and its prospective
training/tuning protocol is followed and disclosed. No accuracy threshold
determines validity: 100\%, 37\%, and random-level accuracy are all retained
outcomes of valid configurations. Low performance can support a conclusion
conditional on the tested model, information, and budget; it does not establish
that a complete algorithm family has failed. An equal-quality parameter
comparison is simply unavailable when the measured qualities differ, rather
than making either run invalid. All failures, configuration trials, retries,
maintenance, and extra training remain charged. A low result is not a software
defect and does not trigger trials until a desired number appears.

\subsection{Five-layer evaluation scope}

\begin{table}[t]
\centering
\small
\begin{tabular}{p{0.07\linewidth}p{0.40\linewidth}p{0.43\linewidth}}
\toprule
Layer & Required scientific coverage & Current evidence boundary\\
\midrule
I & Six-stage acquisition at $N=8,16,32$ & Historical streams in
Table~\ref{tab:growth-parameters}; the frozen fully fused Primary must be
retrained under the five-seed protocol.\\
II & Noncommuting programs, unseen combinations, lengths 16/32/64, dynamic
write/overwrite/read/delay memory & Finite held-out words and length 16 tested;
the complete complex-program and dynamic-memory suite is not completed.\\
III & One Block for reversible, nonreversible, continuous, and variable-length
computation & Current finite reversible operations do not establish this
generality.\\
IV & Speech Commands v2 and Split CIFAR-100 & No competitive real-data result
is claimed in this manuscript checkpoint.\\
V & Applicable architectures, continual-learning controls, and mechanism
ablations & Historical valid measured configurations and working program
controls; prospective bounded family comparisons remain incomplete.\\
\bottomrule
\end{tabular}
\caption{All five layers require direct measurements. Evidence at one layer
does not replace another layer.}
\label{tab:benchmark-layers}
\end{table}

Layer II will distinguish operation-order dependence, including $AB\ne BA$
and commutators, from memorization of an observed word distribution. Programs
are trained at length at most eight and evaluated separately at lengths
16, 32, and 64. A separate key--value stream requires writing, overwriting,
reading, and delayed queries. Every architecture receives terminal task feedback
without teacher intermediate execution traces. Layer III keeps the Block's
core state-transition and learning mechanism fixed across task types; a change
to that core is reported as an architecture variant, rather than silently
counted as generalization.

\subsection{Information, development budgets, and final evaluation}

The synthetic comparison uses the same terminal data, stage order, public
context, and evaluation protocol. Dedicated constructors, teacher operator
tables, hidden state labels, and correct routes are not free priors for any
method. Formation passes and terminal-label accesses beyond minibatch training
are charged explicitly. Equal gradient-update counts do not imply equal total
learning compute when one method additionally solves a full observation table.

The locked synthetic update limits are B1: 600, B2: 2,400, and B3: 9,600
gradient updates per stage, available to every method including Prometheus.
Each main model family has at most 24 development configurations: development
seed 0 selects candidates, the top three are checked with development seed 1,
and the ranking uses a frozen validation criterion. After configuration lock,
the final evaluation uses seeds $0,1,2,3,4$ without test-based tuning. Shared
early checkpoints and resumed run fragments count as one lineage, not multiple
independent seeds. The real-data budgets are 10, 50, and 200 epochs per phase
under the same validation-only selection principle.

Only a localized, reproducible software, data, or training-logic error qualifies
as a bug fix. Optimizer, architecture, curriculum, or initialization changes
are configuration development unless they repair such a demonstrated error.
Any supplemental trajectory outside the registered main budget is disclosed
separately. Reaching the highest budget without acquisition remains an allowed
negative result. Test and OOD measurements do not select configurations,
training duration, or stopping points.

Three comparisons are kept distinct: identical information/protocol, matched
training and total-compute budgets, and equal attained capability quality.
The final analysis requires accuracy--parameter, accuracy--training-cost, and
accuracy--inference-cost frontiers. Different method families may use the
optimizers and schedules appropriate to their original algorithms; forcing one
recipe on all competitors is not a fairness requirement.

\subsection{Fixed real-data protocols and strong competitors}

Speech Commands v2 uses 35 word classes in seven class-incremental phases of
five classes. It uses the official train/validation/test split and common
16-kHz preprocessing. Each model classifies all seen words at inference
without a phase identifier or true class group. Split CIFAR-100 uses ten phases
of ten fine classes, the official 50,000 training and 10,000 test images, and a
fixed class-stratified 10\% training holdout for validation. It uses an expanding
single head and supplies no task ID at inference. Input processing,
augmentation, and data access are matched across architectures.

Both datasets have Memory-0 and Memory-500 conditions, where 500 is the total
number of stored old examples across all classes. Stored features, teacher
outputs, task statistics, and other historical state are additionally recorded;
they are not an exemption from the memory comparison. Main runs start from
random initialization. Pretrained variants are supplemental and give every
method the same pretraining information. These no-ID class-incremental
protocols are distinct from the public-context synthetic benchmark.

The primary competing architecture family is the Transformer: vanilla, replay,
adapters, and replay plus adapters, at two or more prospectively selected scales.
Parameter-matched and quality-matched comparisons are both required.
Continual-learning comparisons include DEN, PackNet, replay, and ER-ACE or an
equally mature class-incremental method. PackNet task-incremental results, when
applicable, are explicitly separated from no-ID class-incremental rankings.
Mamba-2 participates in appropriate sequence/memory tasks, and Titans in tasks
suited to its memory interface. HOPE/Nested Learning participates only with a
verifiable, applicable implementation. An unfaithful reproduction or untested
implementation is marked \texttt{NOT\_VALIDATED}, with its specific fidelity
or implementation gap stated. Low accuracy alone never assigns that status.

\subsection{Measured baseline results at the historical B1 budget}

\begin{table}[t]
\centering
\small
\begin{tabular}{lrrrrrrr}
\toprule
 & & \multicolumn{3}{c}{First task} & \multicolumn{3}{c}{New task at phase 6}\\
Method & Total $P$ & IID & OOD & L16 & IID & OOD & L16\\
\midrule
Transformer-64 &103752&.1836&.1230&.1348&.1152&.1094&.1016\\
Transformer-64+Replay &103752&.1836&.1230&.1348&.1191&.1172&.1055\\
Transformer-96 &229352&.1348&.1094&.1035&.1328&.1406&.1191\\
Transformer-96+Replay &229352&.1348&.1094&.1035&.1133&.1289&.1152\\
\bottomrule
\end{tabular}
\caption{Actual $N=8$, seed-0 historical streams, each with 600 updates per
phase and 3,600 updates overall. ``Phase 6'' is the newly introduced task at
the final phase, not an average over old tasks. These are valid measured B1
configurations under their reported historical protocol, regardless of their
low acquisition. They are not pooled with a future frozen Primary evaluation.}
\label{tab:historical-transformers}
\end{table}

All four runs in Table~\ref{tab:historical-transformers} have zero tasks passing
the dual-interface capability criterion. Their low first-task performance
shows that the result is not explained by forgetting alone. The $N=16$
width-64 Transformer similarly has no passing task in the measured stream.
These are valid budget- and configuration-conditional observations. They do
not by themselves characterize every Transformer configuration or establish an
equal-quality parameter ratio, and they are not new Primary Prometheus results.

A separately recorded anchored Transformer trial uses a 100-update source
prefix and 2,200 additional updates, totaling 2,300 updates. It is a supplemental
configuration trajectory, not a software bug fix or an equal-600-budget result.
It is valid configuration evidence under its own disclosed supplemental protocol.
The fixed development progression stops at its prospective budget. The
validation-selected frozen checkpoint is then tested once, obtaining IID
$0.37109375$, OOD $0.134765625$, and length-16 $0.111328125$. This improves
some short-word behavior without establishing full-task competence; the low
test result is not a trigger for unbounded additional trials.

The first 600-update task runs of the adapted DEN-RNN and linear
self-modifying CMS reconstruction obtain IID $0.115234375$ and
$0.138671875$, respectively. Their parameter counts are 1,800 and 81,224.
The former adapts DEN-style selection, expansion, splitting, and timestamps to
a recurrent terminal-feedback interface; the latter is not an official HOPE
reproduction. These are valid executed adaptations under their reported
600-update protocols; their low scores do not invalidate them. Conclusions
apply to these particular DEN-RNN and CMS variants. Faithful full-original
DEN and HOPE reproduction remains \texttt{NOT\_VALIDATED} because that fidelity
has not been established, not because an accuracy threshold was missed.

For the historical finite-program question, matched conventional controls are available.
At $N=8$, seed 0, ordinary constrained compact learning matches all six
stage-boundary capabilities using 5,310 parameters, versus 5,314 for the
topology model. Repaired freezing uses 6,384 and separate task modules 8,664.
The frozen repair adds a compatible 128-parameter readout when terminal-profile
output multiplicities require it, and later tasks reuse that readout. Its
initial failed version remains recorded. These controls explain the historical
shared/compressed parameter benefit without attributing it to monodromy; they
complement rather than replace the required Transformer comparison.

\subsection{Costs and causal attribution}

\begin{table}[t]
\centering
\small
\begin{tabular}{lrrrrr}
\toprule
 & Weight & Extra-state & Optimizer & Checkpoint & Old guard\\
Policy & bytes & JSON bytes & tensor bytes & bytes & seconds\\
\midrule
Topology compact &21256&208326&42544&341373&47.4276\\
Ordinary constraints &21240&104105&42508&202737&45.6787\\
Frozen + readout &25536&83598&7568&146821&0\\
Private modules &34656&81825&12316&160265&0\\
\bottomrule
\end{tabular}
\caption{Historical final-phase $N=8$, seed-0 resource fields for the
native/coordinate-renderer policies, not the future fully fused Primary.
The guard covers the phase's 600 updates.
These are format-specific fields and overlapping cost scopes, not a disjoint
sum or a controlled wall-time ranking.}
\label{tab:resource-fields}
\end{table}

All stored and frozen parameters count. Reference arrays are contained within
extra-state metadata; cache maps and certificates can also share storage.
Checkpoint sizes include optimizer, RNG, and container state. JSON size is not
minimal sufficient information or peak RAM. Phase wall time includes nested
training, evaluation, and protection scopes, which are not added twice.
Some jobs overlap on the GPU, so their individual wall times do not support a
throughput ranking. One ordinary-control initial-stage report is missing
post-training timing fields after a reload-helper failure; these remain null,
rather than being reconstructed as measured values.

Complete comparison additionally requires development and failed-run costs,
primitive discovery, data accesses, solver and installation work, temporary
dense tensors, replay, cache preparation, old-binding checks, rollback, and
checkpoint reload. Controlled same-hardware training FLOPs/GPU time,
end-to-end inference p50/p95, throughput, and peak memory are not yet complete.
Neither software simulation nor these parameter counts measure hardware energy
or establish a quantum advantage.

The causal matrix removes, separately, primitive reuse, sharing, workspace
growth, old capability protection, graph constraints, and native compilation.
Topology-specific comparisons substitute ordinary constraints for monodromy,
disable cross-class loop learning, disable within-class safe deformations,
and compare ordinary cross-class optimization and discrete search under the
same information and costs. Shared compression, program formation, graph
closure, and generic consistency repair are not credited to monodromy. A
same-mechanism reimplementation tests executable correspondence; it is not by
itself a historical-priority finding or a general impossibility result.

\subsection{Statistical reporting and completion status}

Formal comparisons use all five independent final seeds, mean, standard
deviation, and 95\% confidence intervals; paired-seed differences and their
intervals are primary. Multiple primary tests receive multiplicity correction,
and failed seeds remain included with their actual censoring or failure status.
The principal outcomes are final/average-incremental accuracy, forgetting,
backward transfer, acquisition cost, and total parameters and resources at
equal retained quality. The earlier scale/seed lineages remain historical or
internal ablations; they cannot count as retraining of the newly frozen Primary,
be pooled across architecture versions, or supply its five-seed statistics.

Publication-grade benchmarking remains \texttt{INCOMPLETE}: the fully fused
TODO~2 and TODO~3 must both pass; the complete architecture is then frozen
and actually retrained
across every layer against protocol-valid competitors under the bounded
development procedure. The two fixed real-data benchmarks, causal attribution,
and full resource/statistical evidence remain required. No competitor must
cross an accuracy gate to have its measured result included, and no favorable
historical finite-task result substitutes for the new Primary evaluation.

\paragraph{Evidence locations.}
Original neural streams and their actual checkpoints reside under
\path{.deep-work/2026-10-31T234403-0400-v2/probes/topology-architecture/runs/},
including \path{parameter-n8-seed0-transformer64-full},
\path{parameter-n8-seed0-transformer64-anchored-competence-v1},
\path{parameter-n8-seed0-den32-through0-v2}, and
\path{parameter-n8-seed0-selfmod_cms64-l1-through0-v2}.
The frozen stage-level extract is
\path{.deep-work/2026-10-31T234403-0400-v2/probes/topology-architecture/parameter-boundary-growth-evidence.csv}.
Actual original-source reading and adaptation boundaries for DEN are in
\path{.deep-work/2026-10-31T234403-0400-v2/sources/parameter-den/reading-and-adaptation.md}.
The governing prospective benchmark protocol is \path{ALL.md}, TODO~5, with
the latest explicit human directives controlling architecture identity and
accuracy-independent configuration validity if older wording conflicts.
